\begin{styletable}

Вектор
    & $\vct{x}$, \; $\vct{\theta}_i$
    & \verb|\vct{x}|, \verb|\vct{\theta}_i|
\\

Нулевой вектор
    & $\vctzero$, \; $\vctzero_n$
    & \verb|\vctzero|, \verb|\vctzero_n|
\\

Матрица
    & $\mtr{A}$, \; $\mtr{A}^{-1}$
    & \verb|\mtr{A}|, \verb|\mtr{A}^{-1}|
\\

Нулевая матрица
    & $\mtrzero$, \; $\mtrzero_{m \times n}$
    & \verb|\mtrzero|, \verb|\mtrzero_{m \times n}|
\\

Ранг матрицы
    & $\rank \mtr{A}$
    & \verb|\rank \mtr{A}|
\\

След матрицы
    & $\tr \mtr{A}$
    & \verb|\tr \mtr{A}|
\\

Определитель матрицы
    & $\det \mtr{A}$
    & \verb|\det \mtr{A}|
\\

Размерность пространства
    & $\dim X$
    & \verb|\dim X|
\\

\hline
\end{styletable}

\stylenote{
    Векторы набираются полужирным курсивом строчными буквами ($\vct{x}$, $\vct{\theta}_i$), а матрицы~--- полужирным прямым шрифтом прописными буквами ($\mtr{A}$).
    Для нулевого вектора используется отдельная команда $\vctzero$, поскольку цифры не имеют полужирного курсивного математического начертания.
    Для нулевой матрицы аналогично используется команда $\mtrzero$.
    Компоненты вектора или матрицы являются числами и потому набираются обычным курсивом: $\vct{x} = \tup{x_1, x_2, \ldots, x_n}$.
    Стрелка над буквой для обозначения векторов в книге не используется, однако в рукописном тексте (конспекты, записи на доске), где полужирное начертание невозможно, рекомендуется её применение.
}